% !TeX root = ../Report.tex
\chapter{Conclusion}
\label{sec:conclusion}

This thesis investigated the use of a multi-agent, tool-based architecture for LLM-driven decision support in a satellite ground-segment operational context. The main objective was to address a limitation of conventional LLM-based approaches, namely the difficulty of effectively exploiting large and heterogeneous operational contexts when relevant information is distributed across different sources. To this end, FUNES was designed as an orchestration layer around an underlying language model, introducing explicit mechanisms for information planning, source-specific retrieval, evidence refinement, and response synthesis.

The proposed architecture was evaluated through a controlled benchmark consisting of twenty operational questions divided into three task-complexity categories. The evaluation compared FUNES with Llama 3.3 70B Instruct operating directly on the complete operational context. The results show that FUNES achieved an overall score of $124/138$ ($89.9\%$), compared with $87/138$ ($63.0\%$) for the baseline, corresponding to a difference of 26.9 percentage points. FUNES obtained a higher aggregate score in all three task-complexity categories, with the largest difference observed for easy tasks and progressively smaller differences for medium and hard tasks. The question-level analysis further showed that the advantage of FUNES was particularly pronounced for retrieval-oriented tasks and for cases where relevant evidence had to be selected from heterogeneous operational sources. At the same time, the results included cases in which the baseline achieved equal or higher performance, particularly when the task required detailed operational synthesis or a precise mapping between evidence and requested operational consequences.

These results provide empirical support for the overall design principle underlying FUNES: explicitly controlling the acquisition and organization of operational information can improve the quality of LLM-generated decision-support responses compared with providing a general-purpose LLM with the complete operational context. In particular, the observed improvements are consistent with the hypothesis that source selection and contextual filtering can reduce the impact of irrelevant information and facilitate the subsequent reasoning process. However, the evaluation does not establish that any individual architectural component is responsible for the observed improvement. The absence of an ablation study, the limited number of questions, the use of a single operational scenario and baseline model, and the reliance on automated evaluation limit the extent to which the results can be generalized beyond the considered benchmark.

An important outcome of the evaluation is therefore that information retrieval and orchestration should not be considered a substitute for the reasoning and synthesis capabilities of the underlying LLM. The results for the more complex tasks show that retrieving the correct evidence is not sufficient when the system must transform that evidence into a precise operational judgement, preserve detailed temporal and causal relationships, or formulate a specific operational recommendation. This suggests that future development of FUNES should focus not only on improving information acquisition, but also on requirement-aware synthesis, evidence provenance, uncertainty handling, and iterative validation of the final response.

Overall, FUNES demonstrates the potential of an explicit information-orchestration architecture for integrating LLMs into satellite ground-segment decision-support workflows. Within the scope of the evaluated scenario, the architecture achieved a substantial improvement over the considered full-context baseline while also revealing clear conditions under which this advantage is reduced or absent. Future work should therefore extend the evaluation to larger and more diverse benchmarks, perform controlled ablation studies, investigate different underlying LLMs and reasoning strategies, and validate the system with domain experts in realistic operational environments. These extensions would allow the contribution of the individual architectural mechanisms to be isolated and would determine whether the improvements observed in the current benchmark translate into measurable benefits for real-world operational decision support.